% =====================================================================
%  sec1.tex — 1. The A-Model And The B-Model
%  原文 pp.1–2
% =====================================================================
\section{A-模型与 B-模型}

设 $G$ 为紧李群，$G_{\C}$ 为其复化；又设 $C$ 为紧致的、无边的定向二维流形。我们记 $\YC(G,C)$（在不致混淆时简记为 $\YC(G)$ 或 $\YC$）为 $C$ 上平坦 $G_{\C}$-丛 $E\to C$ 的模空间，模去规范变换。等价地，$\YC(G,C)$ 参数化 $C$ 的基本群到 $G_{\C}$ 的同态\yz{其实，最好把 $\YC(G,C)$ 定义为几何不变量理论（GIT）商：它参数化稳定同态，外加半稳定同态的等价类。这一细化此处与我们无关。见第 6.1 节。}\zhnt{1.1}。

$\YC(G,C)$ 自然地是一个复辛流形，即带有非退化全纯 2-形式的复流形。复结构直接来自 $G_{\C}$ 的复结构；辛形式则记作 $\Omega$，来自 $H^{1}(C,\ad(E))$ 上的相交配对\yz{这一相交配对的定法依赖于 $G$ 的李代数上一个不变二次型的选择。利用 Hitchin 在 Higgs 丛模空间上的 $\C^{*}$ 作用可以证明：我们即将定义的 A-模型与此选择无关，至多差一个自然同构。我们最终要定义的几何 Langlands 对偶同样不依赖这一选择。}\zhnt{1.2}，其中 $\ad(E)$ 是平坦丛 $E$ 的伴随丛。由 $\YC(G,C)$ 是复辛流形这一事实，特别可以推出：它的典范线丛自然地平凡\zhnt{1.3}。

几何 Langlands 对偶所关心的，是与 $\YC(G,C)$ 相联系的一类拓扑场论。其中最基本的是两个：B-模型，其定义方式是把这个流形看作典范丛平凡的复流形；以及 A-模型，其定义方式是把它看作以 $\omega=\operatorname{Im}\Omega$ 为辛形式的实辛流形\yz{$\Omega$ 的通常定义使得 $\operatorname{Im}\Omega$ 在上同调中平凡，而 $\operatorname{Re}\Omega$ 不平凡。$\omega=\operatorname{Im}\Omega$ 上同调平凡这一事实，部分地解释了上一条脚注提到的现象：$\YC$ 的 A-模型在 $\omega$ 的缩放之下不变。}\zhnt{1.4}。

这些就是与最基本形式的几何 Langlands 对偶相关的拓扑场论。不过此外还有一个推广，它与有时被称作量子几何 Langlands 的东西相关。从 A-模型一侧看，推广的可能性是显然的：在定义 A-模型时，我们本可以使用 $\operatorname{Re}\Omega$ 与 $\operatorname{Im}\Omega$ 的更一般的线性组合。不那么显然的是，B-模型作为一个拓扑场论，竟然可以形变到这一族 A-模型之中。这个颇令人意外的事实，从广义复几何的角度看是自然的，见 Hitchin (2003)；Gualtieri (2003) 第 4.6 节正是从这一角度把它作为关于复辛流形的一个一般陈述加以解释的\zhnt{1.5}。在 Kapustin 与 Witten (2007) 第 5.2 与 11.3 节中已经证明：量子几何 Langlands 正是在这一框架中得到自然理解的\zhnt{1.6}。

不过，为使讨论保持简单，下文只关注刚刚描述过的最基本的 B-模型与 A-模型。
